\chapter{Frequent Bowel Movements}

\section*{Definition}
Increased frequency of defecation exceeding three bowel movements per day, often with loose or urgent stools, reflecting altered gastrointestinal motility or secretion.

\section*{Pathophysiology}
Increased stool frequency results from reduced colonic transit time (rapid transit prevents water absorption leading to loose stools), increased secretion (enterotoxins, bile acid malabsorption), or increased smooth muscle contractility (IBS, hyperthyroidism). Acute onset suggests infection (gastroenteritis) while chronic suggests functional IBS, IBD, malabsorption, or medication effect.

\section*{Causes}
\begin{itemize}
  \item Acute gastroenteritis (viral, bacterial, parasitic)
  \item Irritable bowel syndrome (IBS-D)
  \item Inflammatory bowel disease (Crohn disease, ulcerative colitis)
  \item Food intolerances (lactose, gluten, FODMAPs)
  \item Hyperthyroidism
  \item Bile acid malabsorption
  \item Medication side effects (metformin, antibiotics, laxatives, SSRIs)
  \item Anxiety or stress-related
  \item Celiac disease
  \item Colon cancer or carcinoid syndrome (rare)
\end{itemize}

\section*{When to See a Doctor}
See your doctor if:
\begin{itemize}
  \item Severe or worsening symptoms
  \item Frequent bowel movements with blood, weight loss, fever, nocturnal symptoms, or dehydration
  \item Accompanied by other concerning signs
\end{itemize}

\section*{Related Symptoms}
\begin{itemize}
  \item Diarrhea
  \item Abdominal pain
  \item Nausea
  \item Urgency
  \item Bloating
\end{itemize}

\textbf{Important:} If this symptom persists, your doctor will check for serious underlying causes.

\section*{Self-Care / Home Management}
\begin{itemize}
  \item Rest and avoid activities that make it worse.
  \item Maintain adequate hydration and a balanced diet.
  \item Over-the-counter remedies may provide symptomatic relief (consult a pharmacist or doctor).
  \item Monitor symptoms and seek professional advice if they persist or worsen.
  \item Avoid self-medicating with prescription medications without medical guidance.
\end{itemize}

\section*{How Your Doctor Will Evaluate You}
\begin{itemize}
  \item A thorough history and physical examination are the first steps in evaluation.
  \item Laboratory tests (blood work, urinalysis) may be ordered based on clinical suspicion.
  \item Imaging studies (X-ray, ultrasound, CT, MRI) may be indicated when structural pathology is suspected.
  \item Specialist referral is arranged when the diagnosis remains uncertain or requires specific expertise.
\end{itemize}

\section*{Treatment Options}
\begin{itemize}
  \item Treatment targets the underlying cause identified through diagnostic evaluation.
  \item Supportive care and symptom management are provided while awaiting definitive therapy.
  \item Medications, lifestyle modifications, and physical therapies may be prescribed as appropriate.
  \item Follow-up ensures treatment effectiveness and allows timely adjustments to the care plan.
\end{itemize}

\clearpage